% Compile with pdflatex, or paste into Overleaf (overleaf.com).
\documentclass[11pt]{article}
\usepackage[letterpaper, margin=1in, top=0.8in, bottom=0.8in]{geometry}
\usepackage[T1]{fontenc}
\usepackage[utf8]{inputenc}
\usepackage{palatino}
\usepackage{microtype}
\usepackage{amsmath}
\usepackage{amssymb}
\usepackage{titlesec}
\usepackage{enumitem}

\titleformat{\section}{\normalfont\normalsize\bfseries}{\thesection.}{0.5em}{}
\titlespacing*{\section}{0pt}{1em}{0.3em}

\setlength{\parskip}{0.5em}
\setlength{\parindent}{0pt}
\setlist{itemsep=0.15em, topsep=0.2em, partopsep=0pt, leftmargin=1.4em}

\begin{document}

\begin{center}
  {\large\bfseries Epic Games Store Giveaways and Steam Player Engagement}\\[0.4em]
  {\small Nicholas Liu \quad|\quad 14.33 Project Proposal}
\end{center}

\vspace{0.5em}
\hrule
\vspace{0.8em}

\section{Research Question}
What is the causal effect of Epic Games Store free game giveaways on the equivalent game's Steam concurrent player numbers? Does a zero-price substitute on a rival platform \textit{cannibalize} or \textit{promote} engagement on Steam?

Standard theory predicts substitution: rational consumers should prefer the free option, draining Steam's player base. Yet online discussion suggests a counter-intuitive result where Steam counts actually \textit{rise} during Epic giveaways.

\section{Motivation}
\begin{itemize}
  \item \textbf{Platform differentiation.} Steam's moat may be resilient even to zero-price competition.
  \item \textbf{Antitrust policy.} If Epic's subsidies benefit Steam, the platforms function as complements rather than substitutes---or Steam's incumbency is simply too entrenched.
  \item \textbf{Developer strategy.} An Epic giveaway may serve as a marketing event, not a revenue sacrifice.
\end{itemize}

\section{Identification Strategy}
Staggered difference-in-differences (``not-yet-treated'' design): games scheduled for future giveaways serve as controls for currently treated games, avoiding external propensity score matching.

\section{Specification}
\vspace{-0.3em}
\begin{equation}\label{eq:eventstudy}
\ln y_{it} = \alpha_i + \gamma_t + \sum_{\substack{k \in [-12,12] \\ k \neq -1}} \beta_k \cdot \mathbf{1}[t - \mathrm{event}_i = k] + \varepsilon_{it}
\end{equation}
\noindent where $y_{it}$ is average concurrent players for game $i$ in month $t$; $\alpha_i$, $\gamma_t$ are game and month fixed effects; and $\beta_k$ are relative-time dummies ($k=-1$ omitted as baseline).

\section{Hypothesis}
Parallel trends: $\beta_k = 0$ for $k < 0$. Post-treatment: $\beta_k > 0$ indicates complementarity; $\beta_k < 0$ indicates substitution.

\section{Data}
\begin{enumerate}
  \item \textbf{Treatment:} Epic Games Free Giveaway History (2018--2025), Kaggle; $>$500 events.
  \item \textbf{Outcome:} Monthly average/peak concurrent players from SteamCharts.com (scraped).
  \item \textbf{Merge:} Fuzzy string matching on game titles with manual verification.
\end{enumerate}

\end{document}
